\section{Conventions and disc functionals}\label{sec:preliminaries}

Throughout, $X$ is a compact connected K\"ahler manifold of complex
dimension $n\ge1$, with no nonconstant entire curves, and $k$ is a fixed
K\"ahler metric. This section fixes the signs and normalizations used
in the disc estimates. We write $dd^c=i\partial\bar\partial$ and
identify a Hermitian tensor with its real
$(1,1)$-form so that the tensor associated with a real-valued potential
$\phi$ has quadratic form
\[
 \sum_{i,j}\partial_i\partial_{\bar j}\phi\,
 v^i\overline{v^j},
\]
and associated form $\ddc\phi$. In matrix formulas, vectors are columns
and their squared norm is written $v^*hv$. In particular, if the columns
of $B$ are tangent vectors, their Gram matrix is $B^*hB$.

For a K\"ahler form $h$, set
\begin{equation}\label{eq:ricci-convention}
 R_h=-\ddc\log\det h.
\end{equation}
This globally defined form represents $2\pi c_1(T^{1,0}X)$; hence
$[-R_h]=2\pi c_1(K_X)$. Inequalities of real $(1,1)$-forms are understood
as inequalities of their Hermitian tensors. The form $h^n$ is used
without a factor $n!$; determinant ratios are consequently ratios of
these volume forms.

Write $\D=\{z\in\C:|z|<1\}$ and let $\dd S$ denote Euclidean area.
For a holomorphic map $f:\D\to X$, put
\[
 f'(z)=\dd f_z(\partial_z),\qquad d(f)=k(f'(0),f'(0)).
\]
The disc has \emph{finite area} if $\int_\D k(f',f')\,\dd S<\infty$.
Since $X$ is compact, this condition is independent of the smooth
Hermitian metric used to define it. For a smooth closed real
$(1,1)$-form $q$ and a finite-area disc, define
\begin{equation}\label{eq:disc-functionals}
 A_q(f)=\frac1\pi\int_\D q(f',f')\,\dd S,
 \qquad
 P_q(f)=\frac2\pi\int_\D\log\frac1{|z|}\,q(f',f')\,\dd S.
\end{equation}
Both are finite: the weight is bounded away from the center, and the
integrand apart from the weight is smooth near the center. We write
$A_{q,r}(f)$ for the same area integral restricted to $|z|<r$.

For a scalar function on a circle, use the normalized mean
\[
 \langle u\rangle_r=\frac1{2\pi}\int_0^{2\pi}u(re^{i\theta})\,\dd\theta.
\]
The following normalization will be useful repeatedly.

\begin{lemma}[Green identities]\label{pre:green}
If $u$ is smooth on $\D$, then, for $0<r<1$,
\begin{align}
 \langle u\rangle_r-u(0)
 &=\frac2\pi\int_{|z|<r}\log\frac r{|z|}\,
       \partial_z\partial_{\bar z}u\,\dd S,
       \label{eq:green-poisson}\\
 \frac1\pi\int_{|z|<r}\partial_z\partial_{\bar z}u\,\dd S
 &=\frac r2\langle\partial_r u\rangle_r.
       \label{eq:green-area}
\end{align}
If in addition $u$ is continuous on $\overline\D$ and
$\partial_z\partial_{\bar z}u\in L^a(\D)$ for some $a>1$, then
Equation~\eqref{eq:green-poisson} passes to $r=1$.
\end{lemma}

\begin{proof}
The identities are the ordinary planar Green identities with
$\Delta=4\partial_z\partial_{\bar z}$. For the limiting assertion,
$\log(1/|z|)$ belongs to every finite $L^b(\D)$, so H\"older's
inequality and dominated convergence apply to the right-hand side.
Continuity gives convergence of the circle means.
\end{proof}

Sobolev regularity of a continuous map on $\overline\D$ is understood
in finitely many target coordinate charts. The same convention applies
to sections of a pulled-back bundle. When $p>2$, the embedding
$W^{1,p}(\D)\hookrightarrow C^0(\overline\D)$, its local versions, and
the multiplication estimate make $W^{1,p}$ an algebra. Smooth
functions of finitely many Sobolev coefficients are again Sobolev
when their values stay in a compact subset of their domain. We use the
corresponding smooth composition maps between these Banach spaces.
For fiberwise holomorphic functions with uniform radii and Sobolev
coefficient bounds, composition is holomorphic between the complex
Banach spaces in question; the precise application is given in the
variation construction.

The following table collects the notation used throughout the disc
argument. The area and Poisson functionals use the fixed unit disc,
with $A_{q,r}$ denoting the explicitly indicated truncation.
\begin{center}
\begin{tabular}{@{}ll@{}}
\toprule
Notation & Meaning \\
\midrule
$k$ & Fixed background K\"ahler metric on $X$ \\
$h$ & Auxiliary K\"ahler metric \\
$R_h$ & $-dd^c\log\det h$, representing $-2\pi c_1(K_X)$ \\
$d(f)$ & Squared central derivative norm $k(f'(0),f'(0))$ \\
$A_q(f)$ & $\pi^{-1}\int_\D q(f',f')\,\dd S$ \\
$A_{q,r}(f)$ & The same area integral over $|z|<r$ \\
$P_q(f)$ & $2\pi^{-1}\int_\D\log(1/|z|)q(f',f')\,\dd S$ \\
$\langle u\rangle_r$ & Normalized mean of $u$ on $|z|=r$ \\
\bottomrule
\end{tabular}
\end{center}
